\chapter{Analyse et sp\'ecification des besoins}
\label{chap:analyse}

\section{Introduction}
% Annoncer le contenu : acteurs, besoins fonctionnels/non fonctionnels,
% cas d'utilisation.

[\`A compl\'eter]

\section{Identification des acteurs}
% Reprendre les roles deja presents dans l'application (cf.
% docs/cas-utilisation.md) et les decrire en 1-2 phrases chacun.

La plateforme E-FORCE distingue les acteurs suivants~:

\begin{itemize}
    \item \textbf{SuperAdmin} : administre l'ensemble de la plateforme,
    g\`ere les utilisateurs, les processus et les r\`egles m\'etier ;
    \item \textbf{Gestionnaire des processus m\'etier} : cr\'ee et g\`ere les
    processus BPMN et leurs t\^aches, et peut g\'erer les r\`egles m\'etier ;
    \item \textbf{Gestionnaire des r\`egles m\'etier} : cr\'ee, modifie et
    consulte l'historique des r\`egles m\'etier (manuellement ou via le
    module IA/NLP) ;
    \item \textbf{Visiteur / utilisateur non authentifi\'e} : peut
    uniquement se connecter, cr\'eer un compte ou r\'einitialiser son mot
    de passe.
\end{itemize}

\section{Sp\'ecification des besoins fonctionnels}
% Organiser par module. Pour chaque module, lister les fonctionnalites
% sous forme de puces courtes (le detail complet ira dans les cas
% d'utilisation, section suivante).

\subsection{Module Authentification et gestion des utilisateurs}
\begin{itemize}
    \item Authentification via Keycloak (SSO) ;
    \item Inscription avec choix d'un r\^ole m\'etier ;
    \item R\'einitialisation de mot de passe ;
    \item Gestion CRUD des utilisateurs, activation/d\'esactivation des
    comptes, gestion des r\^oles (SuperAdmin uniquement).
\end{itemize}

\subsection{Module Gestion des processus m\'etier (BPMN)}
\begin{itemize}
    \item Cr\'eation, modification, activation/d\'esactivation et
    suppression de processus ;
    \item Gestion des t\^aches d'un processus (formulaires dynamiques,
    r\`egles de transition) ;
    \item G\'en\'eration automatique d'un diagramme BPMN \`a partir des
    t\^aches ;
    \item Visualisation et export du diagramme BPMN.
\end{itemize}

\subsection{Module Moteur de r\`egles m\'etier}
\begin{itemize}
    \item Cr\'eation, modification, activation/d\'esactivation et
    suppression de r\`egles m\'etier (conditions + action + cat\'egorie) ;
    \item G\'en\'eration automatique du code DRL correspondant ;
    \item Gestion du versioning : historique des modifications et
    restauration d'une version ant\'erieure.
\end{itemize}

\subsection{Module IA / NLP}
\begin{itemize}
    \item Saisie d'une phrase en langage naturel (fran\c{c}ais) ;
    \item Conversion automatique en r\`egle m\'etier structur\'ee + code DRL ;
    \item Affichage d'un score de confiance et des ambigu\"it\'es
    \'eventuelles ;
    \item \'Edition manuelle avant validation et enregistrement de la
    r\`egle.
\end{itemize}

\subsection{Module Import / Export de dossiers}
\begin{itemize}
    \item Cr\'eation et consultation des dossiers d'importation et
    d'exportation ;
    \item Filtrage par statut (\texttt{EN\_ATTENTE}, \texttt{EN\_COURS},
    \texttt{VALIDE}, \texttt{REFUSE}) ;
    \item Export des listes filtr\'ees au format Excel (.xlsx) et PDF.
\end{itemize}

\section{Sp\'ecification des besoins non fonctionnels}
% Adapter / completer selon les contraintes reelles du projet.

\begin{itemize}
    \item \textbf{S\'ecurit\'e} : authentification centralis\'ee via
    Keycloak, contr\^ole d'acc\`es bas\'e sur les r\^oles (RBAC), protection
    des routes par garde Angular et v\'erification du jeton JWT \`a chaque
    requ\^ete ;
    \item \textbf{Ergonomie} : interface simple et intuitive, adapt\'ee \`a
    des utilisateurs non techniques (notamment pour la cr\'eation de
    r\`egles via le module NLP) ;
    \item \textbf{Performance} : temps de r\'eponse acceptable pour les
    listes de dossiers, processus et r\`egles ;
    \item \textbf{Portabilit\'e} : application web accessible depuis un
    navigateur, architecture frontend/backend d\'ecoupl\'ee ;
    \item \textbf{Maintenabilit\'e / \'evolutivit\'e} : architecture
    modulaire (composants Angular standalone, services d\'edi\'es par
    domaine).
\end{itemize}

\section{Diagramme de cas d'utilisation global}
% Inserer ici le diagramme de cas d'utilisation global (cf.
% docs/cas-utilisation.md pour le diagramme mermaid existant, a
% reexporter en image ou a redessiner avec un outil UML).

[\`A compl\'eter : diagramme de cas d'utilisation global.]

% \begin{figure}[H]
%     \centering
%     \includegraphics[width=\textwidth]{images/diagramme_cas_utilisation.png}
%     \caption{Diagramme de cas d'utilisation global}
%     \label{fig:cas-utilisation-global}
% \end{figure}

\section{Description des cas d'utilisation principaux}
% Pour chaque cas d'utilisation important, fournir une fiche descriptive
% (acteur, objectif, precondition, scenario nominal, scenarios
% alternatifs, postcondition). Le fichier docs/cas-utilisation.md
% contient deja des fiches pretes a etre reprises/adaptees, par exemple :
%   - Se connecter
%   - Gerer les utilisateurs
%   - Creer un processus / Visualiser et exporter le BPMN
%   - Creer une regle metier (manuellement)
%   - Creer une regle metier via l'IA (NLP)
%   - Consulter l'historique et restaurer une version

[\`A compl\'eter : fiches descriptives des cas d'utilisation principaux.]

\section{Conclusion}
% Resumer le chapitre et faire la transition vers la conception.

[\`A compl\'eter]

\clearpage
